\chapter{Минимальная инвариантная опора спектрального перехода}
\label{ch:v6-spectral-transition-minimal-support-gate}

\section{Проверяемая гипотеза}

Предыдущий гейт заменил буквальную непрерывную нить конечранговым дефектом
оператора перехода
\begin{equation}
 Q_{\mathrm{def}}=|e_0\rangle\langle e_0|\otimes q_0,
 \qquad
 q_0=p_{15}\otimes P_0,
 \qquad
 \rank q_0=15.
\end{equation}
Ранг пятнадцать совпадает с пакетом одного наблюдаемого поколения без
правого нейтрино. Однако совпадение ранга ещё не означает физической
неделимости. Настоящий гейт проверяет нулевую гипотезу: запрещают ли
конечная алгебра, градуировка, Real-структура, заряженный оператор Дирака,
тёплицева граница и локальный перенос любой собственный подпроектор
$r<q_0$.

\section{Контрольный угол не является физической алгеброй}

Связывающий родитель Мориты содержит угол
\begin{equation}
 p_{15}M_{35}(\mathbb C)p_{15}\simeq M_{15}(\mathbb C).
\end{equation}
Его стандартное действие на $\mathbb C^{15}$ неприводимо: совместный
коммутант полного набора матричных единиц равен $\mathbb C I_{15}$.
Если бы этот угол был физической координатной алгеброй, проектор ранга
пятнадцать действительно был бы минимальной ненулевой инвариантной опорой.

Но гейт связывающего родителя уже установил обратную границу: $M_{15}$
является полным контрольным чтением, а не новой gauge-алгеброй. Его
унитарная группа разрешала бы произвольное смешение кварков и лептонов и
была бы значительно больше представленной алгебры Стандартной модели.
Поэтому неприводимость относительно всего $M_{15}$ не является допустимым
доказательством элементарности.

\section{Физическое разложение пакета}

Замороженный носитель равен
\begin{equation}
 H_{15}=Q_L\oplus L_L\oplus u_R\oplus d_R\oplus e_R,
 \qquad
 15=6+2+3+3+1.
 \label{eq:v6-minimal-support-h15}
\end{equation}
Пять слагаемых являются попарно различимыми неприводимыми
gauge-мультиплетами. Ассоциативное замыкание представленной физической
алгебры действует полно внутри каждого неприводимого блока, но не
смешивает блоки с различными квантовыми числами. Поэтому её коммутант на
$H_{15}$ имеет вид
\begin{equation}
 \mathcal A_{\mathrm{phys}}'\big|_{H_{15}}
 \simeq \mathbb C^5.
\end{equation}
На одном gauge-уровне допустимы собственные опоры рангов
\begin{equation}
 6,\ 2,\ 3,\ 3,\ 1.
\end{equation}
Это уже опровергает безусловную физическую примитивность ранга пятнадцать.
Но оператор Дирака может сшить часть этих блоков, поэтому нужен более
сильный тест.

\section{Связность заряженного графа}

На $H_{15}$ проект строго допускает ровно три заряженных ребра
\begin{equation}
 Q_L\xleftrightarrow{\,\widetilde H\,}u_R,
 \qquad
 Q_L\xleftrightarrow{\,H\,}d_R,
 \qquad
 L_L\xleftrightarrow{\,H\,}e_R.
 \label{eq:v6-minimal-support-charged-edges}
\end{equation}
Это обычное чтение конечного графа спектральной тройки: ненулевой блок
оператора Дирака соединяет соответствующие представления, а отсутствие
ребра сохраняет центральный проектор компоненты
\cite{minimal-support-krajewski1998,
minimal-support-chamseddine-connes-marcolli2007}.

Граф \eqref{eq:v6-minimal-support-charged-edges} имеет ровно две связные
компоненты:
\begin{align}
 H_q&=Q_L\oplus u_R\oplus d_R,
 &\dim_{\mathbb C}H_q&=6+3+3=12,\\
 H_\ell&=L_L\oplus e_R,
 &\dim_{\mathbb C}H_\ell&=2+1=3.
\end{align}
Обозначим их проекторы через $P_q$ и $P_\ell$. Тогда
\begin{equation}
 P_q+P_\ell=I_{15},
 \qquad
 P_qP_\ell=0,
 \qquad
 [P_q,\mathcal A_{\mathrm{phys}}]
 =[P_\ell,\mathcal A_{\mathrm{phys}}]=0,
\end{equation}
и для полного заряженного оператора
\begin{equation}
 [P_q,D_{\mathrm{ch}}]=[P_\ell,D_{\mathrm{ch}}]=0.
\end{equation}
Совместный коммутант gauge-действия и всех трёх ненулевых рёбер имеет
комплексную размерность два. Значит, максимально сильное разложение,
которое сохраняется даже при включении всех разрешённых заряженных
связей, равно
\begin{equation}
 H_{15}=H_q\oplus H_\ell,
 \qquad 15=12+3.
 \label{eq:v6-minimal-support-12-plus-3}
\end{equation}
Если одна из амплитуд Юкавы обращается в нуль, число связных компонент
может только увеличиться; склейки кваркового и лептонного пакетов это не
создаёт.

\section{Градуировка, семейная проекция и Real-удвоение}

Хиральная градуировка делит пакет как $8+7$, но оба центральных проектора
$P_q,P_\ell$ с ней коммутируют. Внутри кваркового пакета остаются шесть
левых и шесть правых состояний, внутри лептонного --- два левых и одно
правое. Поэтому градуировка не устраняет разложение
\eqref{eq:v6-minimal-support-12-plus-3}.

Семейный проектор $P_0$ имеет ранг один и действует одинаково на всём
$H_{15}$. Он удаляет семейную кратность, но также не смешивает кварков и
лептонов.

KO6-операция $J$ переводит каждый мультиплет в сопряжённый, сохраняя его
кварковую или лептонную метку. Поэтому Real-инвариантные проекторы имеют
ранги
\begin{equation}
 \rank(P_q\oplus\overline{P_q})=24,
 \qquad
 \rank(P_\ell\oplus\overline{P_\ell})=6,
\end{equation}
и суммируются в полный Real-дефект ранга тридцать. Real-структура удваивает
каждую компоненту, но не сшивает две компоненты друг с другом.

\section{Аддитивность тёплицевой границы}

Коэффициентный проектор раскладывается
\begin{equation}
 q_0=q_q\oplus q_\ell,
 \qquad
 \rank q_q=12,
 \qquad
 \rank q_\ell=3.
\end{equation}
Граничное отображение и фредгольмов индекс аддитивны
\cite{minimal-support-pimsner-voiculescu1980}. Поэтому для положительной
ориентации
\begin{equation}
 -15=-12-3,
\end{equation}
а для обратной ориентации
\begin{equation}
 +15=+12+3.
\end{equation}
Нормированный вес также распадается без остатка:
\begin{equation}
 \frac17=\frac{12}{105}+\frac{3}{105}
 =\frac4{35}+\frac1{35}.
\end{equation}
В каждой Real-компоненте противоположные ориентированные индексы
сокращаются отдельно. Следовательно, ни комплексный boundary-класс, ни
KO6-баланс не требуют удерживать ранги двенадцать и три в одном пакете.

\section{Локальный перенос}

Предыдущая классификация ковариантного переноса дала
\begin{equation}
 W(k)=W_{\mathrm{dir}}(k)\otimes I_E.
\end{equation}
Внутреннее действие скалярно на моритовском носителе. Поэтому
\begin{equation}
 [W,P_q]=[W,P_\ell]=0.
\end{equation}
Перенос предоставляет общий световой конус, но не смешивает кварковую и
лептонную компоненты, не связывает их щели и не защищает ранг пятнадцать
как единый динамический атом.

\begin{theorem}[Составность класса пятнадцать]
Относительно текущей физической алгебры, хиральной градуировки, всех трёх
заряженных рёбер, семейной проекции, KO6-удвоения, тёплицева граничного
класса и локального переноса проектор ранга пятнадцать имеет два
собственных инвариантных подпроектора рангов двенадцать и три. Поэтому
класс пятнадцать является пакетом одного поколения, а не доказанным
элементарным спектральным квантом.
\end{theorem}

\begin{remark}
Ранг три также нельзя автоматически назвать частицей. Это связный
лептонный пакет $L_L\oplus e_R$, содержащий слабый дублет и правое
заряженное состояние. Гейт устанавливает минимальную опору текущего
заряженного графа, но не окончательный атом наблюдаемого спектра.
\end{remark}

\section{Следующий различающий тест}

Следующий гейт
\texttt{version6\_spectral\_transition\_component\_boundary\_gate}
должен отдельно ограничить тёплицев и Real-циклы на $q_q$ и $q_\ell$ и
проверить:
\begin{enumerate}
 \item являются ли оба ограничения настоящими циклами той же конечной
 спектральной тройки, а не только формальными ранговыми классами;
 \item сохраняются ли условие первого порядка и полный физический
 оператор на каждой компоненте;
 \item может ли хотя бы одна компонента получить ненулевую щель без
 независимого секторного параметра.
\end{enumerate}
Нулевая гипотеза следующего гейта: обе компоненты допустимы только как
аддитивные кинематические классы, а их независимое динамическое рождение и
массовые щели текущим родителем не выводятся.

\begin{protocol}
\begin{enumerate}
 \item неприводимость полного контрольного $M_{15}$: \textbf{да};
 \item допустимость $M_{15}$ как физической алгебры: \textbf{нет};
 \item число gauge-блоков $H_{15}$: \textbf{пять};
 \item gauge-коммутант: \textbf{$\mathbb C^5$};
 \item связные компоненты при трёх заряженных рёбрах:
 \textbf{кварковая $12$ и лептонная $3$};
 \item совместный физический коммутант: \textbf{$\mathbb C^2$};
 \item Real-ранги компонент: \textbf{$24$ и $6$};
 \item ориентированные индексы: \textbf{$-12-3=-15$ и
 $+12+3=+15$};
 \item локальный перенос склеивает компоненты: \textbf{нет};
 \item физическая примитивность ранга $15$: \textbf{опровергнута};
 \item ранг $3$ объявлен элементарной частицей: \textbf{нет};
 \item следующий тест: \textbf{граничные и динамические циклы компонент}.
\end{enumerate}
\end{protocol}

Машинный сертификат сохранён в
\path{s2t_v6_spectral_transition_minimal_support_gate_results.json}.

\begin{thebibliography}{9}
\bibitem{minimal-support-krajewski1998}
T.~Krajewski,
\emph{Classification of Finite Spectral Triples},
Journal of Geometry and Physics 28 (1998), 1--30;
arXiv:hep-th/9701081.

\bibitem{minimal-support-chamseddine-connes-marcolli2007}
A.~H.~Chamseddine, A.~Connes, M.~Marcolli,
\emph{Gravity and the Standard Model with Neutrino Mixing},
Advances in Theoretical and Mathematical Physics 11 (2007), 991--1089;
arXiv:hep-th/0610241.

\bibitem{minimal-support-pimsner-voiculescu1980}
M.~Pimsner, D.~Voiculescu,
\emph{Exact Sequences for $K$-Groups and Ext-Groups of Certain
Cross-Product $C^*$-Algebras},
Journal of Operator Theory 4 (1980), 93--118.
\end{thebibliography}